\section{Thermal sensor and image-processing pipeline}
\label{sec:thermal}
\textbf{Simulated, not sensed.} Each part has a heat value (engines and people hot, sun-heated walls warm, ground cool;
the sun adds a solar term by day). Each UAV renders its view of these values into a 320$\times$240 image; a sensor model
turns the ideal image into what an uncooled microbolometer would output, and the computer-vision (CV) pipeline finds
targets in it.
\begin{equation}
y=Q\big(g\odot(h_{\text{PSF}}\ast x)+o+c+n\big),\quad h_{\text{PSF}}=\tfrac1{16}\begin{bmatrix}1&2&1\end{bmatrix}^{\!\top}\!\begin{bmatrix}1&2&1\end{bmatrix},\quad n\sim\mathcal N(0,2^2)
\label{eq:sensor}
\end{equation}
($g$, $o$: per-pixel gain and offset (fixed pattern, partly removed by the shutter NUC), $c$: column stripes, $n$:
temporal noise in DN, $Q$: rounding to 8 bits.) Then (following \cite{gonzalez2018}): Otsu threshold \cite{otsu1979} (with a floor at heat 0.8), despeckle (a hot pixel
needs 2 hot neighbours), connected components by two-pass union-find, blob features, a target filter, and a Kalman
multi-target tracker that fuses both UAVs.
\begin{equation}
G_\sigma(x,y)=\frac{1}{2\pi\sigma^2}e^{-\frac{x^2+y^2}{2\sigma^2}},\quad
S_x=\begin{bmatrix}-1&0&1\\-2&0&2\\-1&0&1\end{bmatrix},\quad
\lVert\nabla I\rVert=\sqrt{(S_x\ast I)^2+(S_y\ast I)^2},\quad s=c\,r^\gamma
\label{eq:kernels}
\end{equation}

\begin{figure}[htbp]
\centering
\raw{lab/cv_pipeline_uav1}{\linewidth}
\caption{The CV pipeline on one UAV frame: raw sensor image, histogram and Otsu threshold, mask, despeckle,
components, overlay.}
\label{fig:cv}
\end{figure}

Sensor noise and warm walls survive the blur but not the threshold and despeckle steps; only compact hot blobs with a
hot core become targets. The CV runs on the CPU at 320$\times$240 on a worker thread, and the PSF blur is separable
(two 3-tap passes instead of one 9-tap pass).
\begin{figure}[htbp]
\centering
\begin{subfigure}{0.49\linewidth}\raw{lab/explorer_imagelab_gauss_gamma}{\linewidth}\caption{Frame Explorer image lab: Gaussian and gamma.}\end{subfigure}\hfill
\begin{subfigure}{0.49\linewidth}\raw{lab/explorer_imagelab_sobel}{\linewidth}\caption{Sobel gradient magnitude.}\end{subfigure}\\[1mm]
{\footnotesize\textit{Zoomed kernel panels of (a) and (b):}}\\[0.5mm]
\feImagelabGaussGammaKernel\\[1mm]
\feImagelabSobelKernel\\[1mm]
\caption{Image-processing stages in the Frame Explorer (\key{`} then \key{P}: \key{K} cycles the kernel, \key{[} \key{]} change $\gamma$).}
\label{fig:imglab-fe}
\end{figure}

\begin{figure}[htbp]
\centering
\begin{subfigure}{\linewidth}\plot{p_thermal_histogram}{\linewidth}\caption{Histogram and CDF before and after equalisation.}\end{subfigure}\\[1mm]
\begin{subfigure}{0.5\linewidth}\centering\plot{p_separable_cost}{\linewidth}\caption{Separable convolution cost.}\end{subfigure}
\caption{Histogram equalisation of a raw thermal frame (a) and the cost of a separable convolution (b).}
\label{fig:imglab}
\end{figure}

The raw thermal frame uses a narrow band of grey levels, because most pixels are cool ground; equalisation spreads them
over the full range, so the CDF becomes a straight line (Figure~\ref{fig:imglab}a). The equalisation is computed
offline for the figure; it is not a switch in the renderer.
\begin{figure}[htbp]
\centering
{\footnotesize\textit{Gaussian ($\sigma=1$, $5\times5$): raw and normalised}}\\[0.5mm]
\wbblock{Thermal_Kernels}{gauss}
{\footnotesize\textit{Thermal patch and the Sobel kernels}}\\[0.5mm]
\wbblock{Thermal_Kernels}{patch}
{\footnotesize\textit{Valid $3\times3$ outputs, gradient magnitude and the thresholded edge map}}\\[0.5mm]
\wbblock{Thermal_Kernels}{out}
\caption{Workbook sheet \code{Thermal\_Kernels}: the normalised Gaussian for the chosen $\sigma$, Sobel kernels, and a
$5\times5$ patch convolved step by step with one output pixel's \code{SUMPRODUCT} shown. \fullsheet{Thermal_Kernels}}
\label{fig:wb-kernels}
\end{figure}

\begin{figure}[htbp]
\centering
\begin{subfigure}{0.49\linewidth}\raw{v5/scene_thermal}{\linewidth}\caption{External thermal view, ironbow palette.}\end{subfigure}\hfill
\begin{subfigure}{0.49\linewidth}\raw{scene/day_thermal_fpv_persons}{\linewidth}\caption{UAV gimbal by day: people classified, not engaged.}\end{subfigure}
\caption{The thermal views. In (b) the ``PERS'' brackets mark warm blobs without a hot engine core: people, never
targeted.}
\label{fig:thermal}
\end{figure}

\textbf{Limitation.} The stages are shown in the lab figures and in the Frame Explorer's image lab (kernel and gamma
controls), but there is no single live strip of all stages with threshold and equalisation sliders.
